\section{The group \texorpdfstring{$SU(2)$}{SU(2)} and its Lie algebra \texorpdfstring{$\mathfrak{su}(2)$}{su(2)}}\label{quat:sec:su2}
This section follows Frankel's Section~19.1b (pp.~493--495) in the
notation of \S\ref{quat:sec:notation}.

\subsection{The group}
As in the Lie-groups part,
\begin{equation}\label{quat:SU2-def}
 SU(2)=\{U\in GL(2,\mathbb C):UU^\dagger=1,\ \det U=1\}.
\end{equation}
Dropping $\det U=1$ gives the unitary group $U(2)$. For square matrices
$UU^\dagger=1$ and $U^\dagger U=1$ are equivalent, and they say that $U$
preserves the Hermitian inner product $\langle z,w\rangle=z^\dagger w$ on
$\mathbb C^2$: $(Uz)^\dagger(Uw)=z^\dagger U^\dagger Uw=z^\dagger w$.

Every element has a standard form. Write
$U=\begin{bmatrix}a&b\\c&d\end{bmatrix}$. Since $\det U=1$, the inverse is
$U^{-1}=\begin{bmatrix}d&-b\\-c&a\end{bmatrix}$; since $UU^\dagger=1$,
the inverse is also $U^\dagger=\begin{bmatrix}\bar a&\bar c\\
\bar b&\bar d\end{bmatrix}$. Equating entries gives $d=\bar a$ and
$-b=\bar c$, that is $b=-\bar c$. Then $\det U=a\bar a-(-\bar c)c
=|a|^2+|c|^2$. Writing $\alpha=a$, $\beta=c$,
\begin{equation}\label{quat:su2-general}
 SU(2)=\Bigl\{\begin{bmatrix}\alpha&-\bar\beta\\ \beta&\bar\alpha\end{bmatrix}
 :\alpha,\beta\in\mathbb C,\ |\alpha|^2+|\beta|^2=1\Bigr\}.
\end{equation}
Conversely, every matrix of this form satisfies $UU^\dagger=1$ and
$\det U=1$, by direct multiplication.

\emph{Dimension.} $GL(2,\mathbb C)$ has four complex entries, so real
dimension $8$. The matrix $U^\dagger U$ is Hermitian for every $U$, so
$U^\dagger U=1$ imposes only four real conditions: the two real diagonal
entries and the real and imaginary parts of one off-diagonal entry. Hence
$\dim U(2)=8-4=4$. On $U(2)$,
$|\det U|^2=\det(U^\dagger)\det U=\det(U^\dagger U)=1$, so
$\det U=e^{i\phi}$ for some real $\phi$, and $\det U=1$ is one further real
condition:
\begin{equation}\label{quat:dim-SU2}
 \dim SU(2)=4-1=3 .
\end{equation}
Equation~\eqref{quat:su2-general} agrees: $(\alpha,\beta)\in\mathbb C^2$
has four real parameters and $|\alpha|^2+|\beta|^2=1$ is one real equation.

\subsection{The Lie algebra}
Let $U(t)$ be a curve in $SU(2)$ with $U(0)=1$ and $X=\frac{dU}{dt}(0)$.
Differentiating $U(t)U(t)^\dagger=1$ by the product rule at $t=0$ gives
$X+X^\dagger=0$. For the determinant use
$\det\exp(tX)=e^{t\operatorname{tr}X}$: by
Theorem~\ref{dgdt:Thm:5.1YKos}(ii), $\exp(tX)\in SU(2)$ for all $t$, so
$e^{t\operatorname{tr}X}=1$ for all $t$; differentiating at $t=0$ gives
$\operatorname{tr}X=0$. Conversely, if $X^\dagger=-X$ and
$\operatorname{tr}X=0$, then $\exp(tX)^\dagger\exp(tX)=\exp(-tX)\exp(tX)=1$
and $\det\exp(tX)=e^0=1$. So
\begin{equation}\label{quat:su2-def}
 \mathfrak{su}(2)=\{X\in\mathfrak{gl}(2,\mathbb C):X+X^\dagger=0,\
 \operatorname{tr}X=0\},
\end{equation}
the definition recorded in \S\ref{dgdt:sec:KS-su2-bases}. Frankel writes
$\mathfrak{su}(2)=\mathfrak g$ and calls its elements skew-Hermitian
matrices with trace $0$.

\emph{Dimension.} $X^\dagger=-X$ makes the diagonal entries purely
imaginary (two real parameters) and fixes $X_{21}=-\overline{X_{12}}$ (two
real parameters in $X_{12}$). Then $\operatorname{tr}X$ is purely imaginary,
and $\operatorname{tr}X=0$ is one real condition. Every
$X\in\mathfrak{su}(2)$ is therefore
\begin{equation}\label{quat:su2-general-X}
 X=\begin{bmatrix}ia&w\\-\bar w&-ia\end{bmatrix},\qquad a\in\mathbb R,\
 w\in\mathbb C,\qquad \dim\mathfrak{su}(2)=4-1=3=\dim SU(2).
\end{equation}
$\mathfrak{su}(2)$ is a vector space over $\mathbb R$, not over
$\mathbb C$: if $X\in\mathfrak{su}(2)$ is nonzero then
$(iX)^\dagger=-iX^\dagger=iX$, so $iX$ is Hermitian and not in
$\mathfrak{su}(2)$.

\emph{Hermitian generators.} Multiplication by $i$ is therefore a real
linear isomorphism from $\mathfrak{su}(2)$ onto the traceless Hermitian
matrices, which Frankel calls $i\mathfrak g$ and we write
$i\,\mathfrak{su}(2)$. Physics texts usually work in
$i\,\mathfrak{su}(2)$ (Hermitian generators, observables) and write group
elements as $\exp(-iH)$ with $H$ Hermitian; Frankel follows this with
$e^{th/i}$.

\emph{Bases.} The Pauli matrices
$\sigma_1=\begin{bmatrix}0&1\\1&0\end{bmatrix}$,
$\sigma_2=\begin{bmatrix}0&-i\\i&0\end{bmatrix}$,
$\sigma_3=\begin{bmatrix}1&0\\0&-1\end{bmatrix}$ are Hermitian and
traceless. The products
\begin{equation}\label{quat:minus-i-sigma}
 -i\sigma_1=\begin{bmatrix}0&-i\\-i&0\end{bmatrix},\qquad
 -i\sigma_2=\begin{bmatrix}0&-1\\1&0\end{bmatrix},\qquad
 -i\sigma_3=\begin{bmatrix}-i&0\\0&i\end{bmatrix}
\end{equation}
give $x^k(-i\sigma_k)=\begin{bmatrix}-ix^3&-x^2-ix^1\\x^2-ix^1&ix^3
\end{bmatrix}$. This equals the $X$ of Eq.~\eqref{quat:su2-general-X} exactly
when $x^3=-a$, $x^2=-\operatorname{Re}w$, $x^1=-\operatorname{Im}w$, so
every $X\in\mathfrak{su}(2)$ is $x^k(-i\sigma_k)$ for unique real $x^k$.
Thus
\begin{equation}\label{quat:su2-bases}
 \{-i\sigma_k\}=\{\sigma_k/i\}\ \text{is a basis of }\mathfrak{su}(2),
 \qquad \{\sigma_k\}\ \text{is a basis of } i\,\mathfrak{su}(2).
\end{equation}
These are Frankel's bases. Holm uses the same basis $-i\sigma_k$ of
$\mathfrak{su}(2)$ (Section~3.2.3, p.~102). Kosmann-Schwarzbach's basis
$\xi_k$ of \S\ref{dgdt:sec:KS-su2-bases} is related to them by
$\sigma_1=-2i\xi_1$, $\sigma_2=2i\xi_2$, $\sigma_3=-2i\xi_3$
(\cite{dgdt:YKosmann-Schwarzbach2010}, Ch.~5 \S1.1, p.~71), so
$\xi_k=\pm\frac12(-i\sigma_k)$ with the signs $(-,+,-)$; the meaning of the
signs is Eq.~\eqref{quat:KS-vs-Frankel}.

\subsection{Vectors as traceless Hermitian matrices}
Define the real linear map
\begin{equation}\label{quat:Sigma}
 \Sigma\colon\mathbb R^3\to i\,\mathfrak{su}(2),\qquad
 \Sigma(\mathbf x)=x^k\sigma_k
 =\begin{bmatrix}z&x-iy\\x+iy&-z\end{bmatrix},
\end{equation}
for $\mathbf x=(x,y,z)$; the matrix is
$x\sigma_1+y\sigma_2+z\sigma_3$ added entry by entry, and
$\Sigma(\mathbf e_k)=\sigma_k$. This is Frankel's Eq.~(19.7), where
$\Sigma(\mathbf x)$ is written $\mathbf x_*=\mathbf x\bullet\boldsymbol\sigma$.
Holm writes the same matrix as $\mathbf r\cdot\boldsymbol\sigma$
(Eq.~(3.2.13), p.~101). Multiplying by $-i$ lands in $\mathfrak{su}(2)$
instead: Holm's tilde map (Eq.~(3.2.14), p.~102) is
$\widetilde{\mathbf q}=-i\,\mathbf q\cdot\boldsymbol\sigma=-i\Sigma(\mathbf q)$,
which is our pure quaternion $q^ae_a$ under Eq.~\eqref{quat:dictionary}.
Kosmann-Schwarzbach identifies $\mathbb R^3$ with $\mathfrak{su}(2)$ by
$a\mapsto X_a=a^1\mathcal I+a^2\mathcal J+a^3\mathcal K=2a^k\xi_k$
(Ch.~5 \S2.2, p.~77), which differs from $-i\Sigma$ in two signs; see
\S\ref{quat:sec:KS-phi}.

To invert $\Sigma$, use Frankel's anticommutator
$\sigma_j\sigma_k+\sigma_k\sigma_j=2\delta_{jk}I$ (Eq.~(19.18)). Taking
the trace, $2\operatorname{tr}(\sigma_j\sigma_k)=2\delta_{jk}\operatorname{tr}I
=4\delta_{jk}$, so $\operatorname{tr}(\sigma_j\sigma_k)=2\delta_{jk}$. Then
\begin{equation}\label{quat:Sigma-inverse}
 \operatorname{tr}\bigl(\Sigma(\mathbf x)\sigma_j\bigr)
 =x^k\operatorname{tr}(\sigma_k\sigma_j)=2x^j,\qquad
 x^j=\tfrac12\operatorname{tr}\bigl(\Sigma(\mathbf x)\sigma_j\bigr),
\end{equation}
Frankel's Eq.~(19.8). With his real scalar product
$\langle h,h'\rangle=\operatorname{tr}(hh')$ on $i\,\mathfrak{su}(2)$
(Eq.~(19.9)),
\begin{equation}\label{quat:Sigma-metric}
 \langle\Sigma(\mathbf x),\Sigma(\mathbf y)\rangle
 =x^jy^k\operatorname{tr}(\sigma_j\sigma_k)=2\,\mathbf x\cdot\mathbf y ,
\end{equation}
so the Pauli matrices are orthogonal, each of length $\sqrt2$. The same
symmetrization as in Eq.~\eqref{quat:sigma-square} gives
$\Sigma(\mathbf x)^2=x^jx^k\sigma_j\sigma_k=|\mathbf x|^2I$, and from the
matrix in Eq.~\eqref{quat:Sigma},
$\det\Sigma(\mathbf x)=-z^2-(x-iy)(x+iy)=-|\mathbf x|^2$.

\subsection{The adjoint representation}\label{quat:sec:adjoint}
Frankel writes ``Recall that every Lie group $G$ acts on its Lie algebra
$\mathfrak g$ by the adjoint action $\mathrm{Ad}\colon G\to Gl(\mathfrak g)$,
$\mathrm{Ad}(g)(X)=gXg^{-1}$'' (p.~494). Two words need definitions.
An \emph{action} of a group $G$ on a set $M$ is a map
$G\times M\to M$, $(g,x)\mapsto g\cdot x$, with $1\cdot x=x$ and
$(gh)\cdot x=g\cdot(h\cdot x)$. When $M=E$ is a vector space and each map
$x\mapsto g\cdot x$ is linear, the action is the same thing as a group
homomorphism $\rho\colon G\to GL(E)$, $\rho(g)x=g\cdot x$: a
\emph{representation} $(E,\rho)$ of $G$, in Kosmann-Schwarzbach's notation
(\S\ref{dgdt:sec:lie-groups}). So the claim is that $\mathrm{Ad}$ is a
representation of $G$ on its own Lie algebra. It is short enough to prove.

\begin{proposition}[adjoint representation]\label{quat:prop:Ad}
Let $G\subseteq GL(n,\mathbb C)$ be a matrix Lie group with Lie algebra
$\mathfrak g$, and for $g\in G$ put $\mathrm{Ad}_gY=gYg^{-1}$. Then
\begin{enumerate}
 \item[(i)] $\mathrm{Ad}_gY\in\mathfrak g$ for every $Y\in\mathfrak g$;
 \item[(ii)] $\mathrm{Ad}_g\colon\mathfrak g\to\mathfrak g$ is linear and
  invertible, with inverse $\mathrm{Ad}_{g^{-1}}$;
 \item[(iii)] $\mathrm{Ad}_1=\mathrm{id}$ and
  $\mathrm{Ad}_{gh}=\mathrm{Ad}_g\circ\mathrm{Ad}_h$, so
  $\mathrm{Ad}\colon G\to GL(\mathfrak g)$ is a representation;
 \item[(iv)] $\mathrm{Ad}_g[X,Y]=[\mathrm{Ad}_gX,\mathrm{Ad}_gY]$, so each
  $\mathrm{Ad}_g$ is a Lie algebra automorphism;
 \item[(v)] $\mathrm{Ad}_g$ is the derivative at $1$ of the conjugation
  $\mathcal C_g(h)=ghg^{-1}$, and
  $\frac{d}{dt}\mathrm{Ad}_{\exp(tX)}Y\big|_{t=0}=[X,Y]=\mathrm{ad}_XY$.
\end{enumerate}
\end{proposition}
\begin{proof}
(i) In $(gYg^{-1})^k$ every inner $g^{-1}g$ cancels, for example
$(gYg^{-1})^2=gYg^{-1}gYg^{-1}=gY^2g^{-1}$. Summing the exponential series,
\begin{equation}\label{quat:conj-exp}
 \exp\bigl(t\,gYg^{-1}\bigr)=\sum_{k=0}^{\infty}\frac{t^k}{k!}gY^kg^{-1}
 =g\exp(tY)g^{-1}=\mathcal C_g\bigl(\exp(tY)\bigr).
\end{equation}
The right side is a product of three elements of $G$, so it lies in $G$
for all $t$, and Theorem~\ref{dgdt:Thm:5.1YKos}(ii) gives
$gYg^{-1}\in\mathfrak g$.

(ii) Matrix multiplication is linear in $Y$. Also
$\mathrm{Ad}_{g^{-1}}\mathrm{Ad}_gY=g^{-1}gYg^{-1}g=Y$, and likewise in the
other order.

(iii) $\mathrm{Ad}_1Y=Y$, and
$\mathrm{Ad}_{gh}Y=ghY(gh)^{-1}=g(hYh^{-1})g^{-1}=\mathrm{Ad}_g(\mathrm{Ad}_hY)$,
using $(gh)^{-1}=h^{-1}g^{-1}$.

(iv) Insert $g^{-1}g=1$ between the factors:
$g(XY-YX)g^{-1}=gXg^{-1}gYg^{-1}-gYg^{-1}gXg^{-1}
=[\mathrm{Ad}_gX,\mathrm{Ad}_gY]$.

(v) By Eq.~\eqref{quat:conj-exp}, $\mathcal C_g$ carries the curve
$\exp(tY)$, whose velocity at $t=0$ is $Y$, to the curve
$\exp(t\,gYg^{-1})$, whose velocity at $t=0$ is $gYg^{-1}$. So the derivative
of $\mathcal C_g$ at $1$ sends $Y$ to $\mathrm{Ad}_gY$. For the last
statement, differentiate $\exp(tX)Y\exp(-tX)$ at $t=0$ by the product rule:
$XY+Y(-X)=[X,Y]$.
\end{proof}
Kosmann-Schwarzbach takes (v) as the definition, $\mathrm{Ad}_g=D_1\mathcal
C_g$, and derives $\mathrm{Ad}_gY=gYg^{-1}$
(\S\ref{dgdt:sec:adjoint-representation}); Frankel takes the formula as the
definition. By (v) the two agree.

\begin{center}
\small
\begin{tabular}{llll}
 & Frankel & Kosmann-Schwarzbach & ours\\\hline
 conjugation & --- & $\mathcal C_g$ & $\mathcal C_g$\\
 adjoint action & $\mathrm{Ad}(g)(X)$ & $\mathrm{Ad}_g=D_1\mathcal C_g$ &
   $\mathrm{Ad}_gX$\\
 its target & $Gl(\mathfrak g)$ & $GL(\mathfrak g)$ & $GL(\mathfrak g)$\\
 representation & ``by $3\times3$ matrices'' & $(E,\rho)$,
   $\rho\colon G\to GL(E)$ & $(E,\rho)$\\
 derivative of $\mathrm{Ad}$ & $\mathrm{Ad}_*$ & $\mathrm{ad}_X$ &
   $\mathrm{ad}_X$, $d\Phi$
\end{tabular}
\end{center}

For $G=SU(2)$, Frankel lets $SU(2)$ act on $i\,\mathfrak{su}(2)$ instead of
$\mathfrak{su}(2)$ and still calls the action $\mathrm{Ad}$ (his
Eq.~(19.10)). This is the same action: $\mathrm{Ad}_u$ is complex linear on
$2\times2$ matrices, so $\mathrm{Ad}_u(iX)=i\,\mathrm{Ad}_u(X)$. Directly,
for Hermitian traceless $h$ and $u^{-1}=u^\dagger$,
$(uhu^{-1})^\dagger=(uhu^\dagger)^\dagger=uh^\dagger u^\dagger=uhu^{-1}$ and
$\operatorname{tr}(uhu^{-1})=\operatorname{tr}(hu^{-1}u)=\operatorname{tr}h=0$.

\subsection{The homomorphism \texorpdfstring{$\Phi$}{Phi}}
Since $u\Sigma(\mathbf x)u^{-1}$ is traceless Hermitian, it is
$\Sigma$ of exactly one vector. Define $\Phi(u)\mathbf x$ to be that vector:
\begin{equation}\label{quat:Phi}
 \Sigma\bigl(\Phi(u)\mathbf x\bigr)=u\,\Sigma(\mathbf x)\,u^{-1}
 =\mathrm{Ad}_u\Sigma(\mathbf x),
 \qquad u\in SU(2),\ \mathbf x\in\mathbb R^3 .
\end{equation}
$\Phi(u)$ is linear in $\mathbf x$, and by Eq.~\eqref{quat:Sigma-inverse}
its matrix entries are
\begin{equation}\label{quat:Phi-entries}
 \Phi(u)_{jk}=\tfrac12\operatorname{tr}\bigl(\sigma_j\,u\sigma_ku^{-1}\bigr).
\end{equation}
Equation~\eqref{quat:Phi} says that this square commutes:
\begin{equation}
\begin{tikzcd}
 \mathbb R^3 \arrow[r,"\Phi(u)"] \arrow[d,"\Sigma"'] &
 \mathbb R^3 \arrow[d,"\Sigma"]\\
 i\,\mathfrak{su}(2) \arrow[r,"\mathrm{Ad}_u"'] & i\,\mathfrak{su}(2)
\end{tikzcd}
\end{equation}
Frankel writes both $\Phi(u)$ and $\mathrm{Ad}_u$ as $\mathrm{Ad}(u)$, and
writes $\mathrm{Ad}\,u\,\mathbf x\sim\mathrm{Ad}\,u\,\mathbf x_*$ for
Eq.~\eqref{quat:Phi}: the $\sim$ means ``corresponds under $\Sigma$''.
Because $\Sigma(\mathbf x)/i\in\mathfrak{su}(2)$ and
$\mathrm{Ad}_u(\Sigma(\mathbf x)/i)=\mathrm{Ad}_u(\Sigma(\mathbf x))/i$,
$\Phi(u)$ is also the matrix of $\mathrm{Ad}_u$ on $\mathfrak{su}(2)$ in
the basis $\sigma_k/i$. So $\Phi$ is the adjoint representation of $SU(2)$
written as $3\times3$ real matrices, which is why Frankel calls it
$\mathrm{Ad}$.

$\Phi$ is a homomorphism. Apply Eq.~\eqref{quat:Phi} twice:
\begin{equation}
 \Sigma\bigl(\Phi(uu')\mathbf x\bigr)=uu'\Sigma(\mathbf x)(uu')^{-1}
 =u\,\Sigma\bigl(\Phi(u')\mathbf x\bigr)u^{-1}
 =\Sigma\bigl(\Phi(u)\Phi(u')\mathbf x\bigr),
\end{equation}
and $\Sigma$ is injective, so $\Phi(uu')=\Phi(u)\Phi(u')$.

$\Phi(u)$ is orthogonal. By Eq.~\eqref{quat:Sigma-metric} and cyclicity of
the trace,
\begin{align}
 2|\Phi(u)\mathbf x|^2
 &=\operatorname{tr}\Bigl(\Sigma\bigl(\Phi(u)\mathbf x\bigr)^2\Bigr)
  =\operatorname{tr}\bigl(u\Sigma(\mathbf x)u^{-1}u\Sigma(\mathbf x)u^{-1}\bigr)\notag\\
 &=\operatorname{tr}\bigl(u\Sigma(\mathbf x)^2u^{-1}\bigr)
  =\operatorname{tr}\bigl(\Sigma(\mathbf x)^2\bigr)=2|\mathbf x|^2 .
\end{align}
A linear map that preserves lengths preserves inner products, since
$2\mathbf x\cdot\mathbf y=|\mathbf x+\mathbf y|^2-|\mathbf x|^2-|\mathbf y|^2$.
So $\Phi(u)^T\Phi(u)=I$ and $\det\Phi(u)=\pm1$. That the sign is $+1$
needs the topology of $SU(2)$.

\subsection{What ``yields the standard representation of \texorpdfstring{$SO(3)$}{SO(3)}'' means}\label{quat:sec:standard-rep}
Frankel announces (p.~493): ``the adjoint representation
$\mathrm{Ad}(g)Y=gYg^{-1}$ of $SU(2)$ on its $3$-dimensional Lie algebra
$\mathfrak{su}(2)$ yields (see Theorem (19.12)) the standard representation
of $SO(3)$ on $\mathbb R^3$.'' Here are the definitions involved.

The \emph{standard} (defining) representation of a matrix group
$G\subseteq GL(n,\mathbb R)$ is the inclusion: $g\in G$ acts on
$\mathbf x\in\mathbb R^n$ by matrix multiplication, $\mathbf x\mapsto g\mathbf x$.
For $SO(3)$ this is the rotation of vectors in $\mathbb R^3$. Composing it
with $\Phi$ gives a representation $(\mathbb R^3,\Phi)$ of $SU(2)$.

For representations $(E_1,\rho_1)$ and $(E_2,\rho_2)$ of a group $G$, an
\emph{intertwining operator} is a linear map $T\colon E_1\to E_2$ with
$\rho_2(g)\,T=T\,\rho_1(g)$ for all $g\in G$. Following Kosmann-Schwarzbach's
Definition~1.12, $\rho_1$ and $\rho_2$ are \emph{equivalent} if there is a
bijective intertwining operator for them. Equivalent representations are the
same representation written in two coordinate systems: $T$ translates one
into the other.

\begin{proposition}\label{quat:prop:standard}
The map $\Sigma\colon\mathbb R^3\to i\,\mathfrak{su}(2)$ is a bijective
intertwining operator from $(\mathbb R^3,\Phi)$ to
$(i\,\mathfrak{su}(2),\mathrm{Ad})$. Likewise
$-i\Sigma\colon\mathbb R^3\to\mathfrak{su}(2)$ is one from $(\mathbb R^3,\Phi)$
to $(\mathfrak{su}(2),\mathrm{Ad})$. Hence the adjoint representation of
$SU(2)$ is equivalent to the standard representation of $SO(3)$ composed with
$\Phi$. By Theorem~\ref{quat:thm:covering}, every rotation occurs as some
$\Phi(u)$.
\end{proposition}
\begin{proof}
Eq.~\eqref{quat:Phi} reads $\mathrm{Ad}_u\circ\Sigma=\Sigma\circ\Phi(u)$
for every $u\in SU(2)$, which is the intertwining condition. $\Sigma$ is
linear, and it is bijective because $\{\sigma_k\}$ is a basis of
$i\,\mathfrak{su}(2)$ (Eq.~\eqref{quat:su2-bases}). Multiplying by $-i$
commutes with $\mathrm{Ad}_u$ and maps $i\,\mathfrak{su}(2)$ bijectively
onto $\mathfrak{su}(2)$. The last sentence is Theorem~\ref{quat:thm:covering}.
\end{proof}
So ``$\mathrm{Ad}$ yields the standard representation'' combines two
statements. In the coordinates $\Sigma$, the operators $\mathrm{Ad}_u$
\emph{are} rotation matrices, by Proposition~\ref{quat:prop:standard} and
$\det\Phi(u)=1$. And all rotations are obtained, by
Theorem~\ref{quat:thm:covering}.

\emph{What suggests this before any computation.} Four facts, each already
available, point to it.
\begin{enumerate}
 \item Dimension: $\dim\mathfrak{su}(2)=3$ (Eq.~\eqref{quat:su2-general-X}),
  so $\mathrm{Ad}_u$ is a $3\times3$ real matrix.
 \item Invariant inner product: $\langle h,h'\rangle=\operatorname{tr}(hh')$
  is positive definite on $i\,\mathfrak{su}(2)$ (Eq.~\eqref{quat:Sigma-metric})
  and $\mathrm{Ad}_u$ preserves it, by cyclicity of the trace. A linear map
  preserving a positive-definite inner product is orthogonal in an
  orthonormal basis, so $\mathrm{Ad}_u\in O(3)$.
 \item Matching brackets: the basis $\sigma_a/(2i)$ of $\mathfrak{su}(2)$
  satisfies $[\sigma_a/2i,\sigma_b/2i]=\epsilon_{abc}\,\sigma_c/2i$
  (Eq.~\eqref{quat:Fa-brackets}), exactly the relations
  $[E_a,E_b]=\epsilon_{abc}E_c$ of $\mathfrak{so}(3)$
  (Eq.~\eqref{quat:so3-brackets}). These are Frankel's ``remarks after
  (19.6)''. The Lie algebras are isomorphic, and by
  Proposition~\ref{quat:prop:Ad}(v) the derivative of $\mathrm{Ad}$ is
  $\mathrm{ad}$, which carries $\sigma_a/2i$ to $E_a$
  (Eq.~\eqref{quat:frankel-19-14}).
 \item Connectedness: $SU(2)$ is connected (\S\ref{quat:sec:su2-sphere}),
  so $\det\mathrm{Ad}_u$ cannot jump from $+1$ to $-1$.
\end{enumerate}
Frankel adds a fifth, which he does not prove here (p.~493): $SU(2)$ is
simply connected, and Lie theory then lifts the Lie algebra isomorphism
$\mathfrak{su}(2)\cong\mathfrak{so}(3)$ to a group homomorphism
$SU(2)\to SO(3)$. In these notes $\Phi$ is constructed directly instead, by
Eq.~\eqref{quat:Phi}.

\section{\texorpdfstring{$SU(2)$}{SU(2)} is the 3-sphere}\label{quat:sec:su2-sphere}
This section follows Frankel's Section~19.1c (p.~495). Let
\begin{equation}
 S^3=\{(z_1,z_2)^T\in\mathbb C^2:|z_1|^2+|z_2|^2=1\}\subset\mathbb C^2\cong\mathbb R^4 .
\end{equation}
A unitary matrix preserves $|z|^2=z^\dagger z$, so each $u\in SU(2)$ maps
$S^3$ to itself. This is what Frankel means by ``$SU(2)\colon S^3\to S^3$;
this is the meaning of being unitary''.

\subsection{The action is transitive}
A group $G$ \emph{acts transitively} on a set $M$ if for every pair of
points $x,y\in M$ some $g\in G$ has $gx=y$ (Frankel, Section~17.2a,
p.~456): one can move from any point to any other. It is enough to reach
every point from one fixed point $x_0$. If $x=gx_0$ and $y=g'x_0$, then
$g'g^{-1}x=y$.

Take $x_0=(1,0)^T$ and any $(z_1,z_2)^T\in S^3$. The matrix
\begin{equation}\label{quat:frankel-19-11}
 u=\begin{bmatrix}z_1&-\bar z_2\\z_2&\bar z_1\end{bmatrix}
\end{equation}
is Eq.~\eqref{quat:su2-general} with $\alpha=z_1$, $\beta=z_2$, so
$u\in SU(2)$. Also $u(1,0)^T$ is the first column of $u$, which is
$(z_1,z_2)^T$. This is Frankel's Eq.~(19.11). So $SU(2)$ acts transitively
on $S^3$.

\subsection{From the orbit to a homeomorphism}
Frankel's Fundamental Principle~(17.10), Section~17.2a, p.~457, states:
\emph{if $G$ acts transitively on $M$, $x_0\in M$, and
$H=\{g\in G:gx_0=x_0\}$ is the stability (isotropy) subgroup of $x_0$, then
$gH\mapsto gx_0$ is a one-to-one correspondence between the left cosets
$G/H$ and $M$.} It comes from the orbit map $g\mapsto gx_0$. This map is
onto by transitivity. Two elements have the same image exactly when they lie
in one coset: $gx_0=g'x_0\iff g^{-1}g'x_0=x_0\iff g^{-1}g'\in H\iff
gH=g'H$. Theorem~(17.11), p.~458, adds that for a closed subgroup $H$ of a
Lie group $G$, $G/H$ is a manifold of dimension $\dim G-\dim H$.

For $x_0=(1,0)^T$, write $u$ in the form~\eqref{quat:su2-general}. Then
$u(1,0)^T=(\alpha,\beta)^T$, and $u(1,0)^T=(1,0)^T$ forces $\alpha=1$,
$\beta=0$, hence
\begin{equation}
 u=\begin{bmatrix}1&-\bar0\\0&\bar1\end{bmatrix}=I,\qquad H=\{I\}.
\end{equation}
The stability subgroup is trivial because Eq.~\eqref{quat:su2-general}
determines the whole matrix from its first column. Each coset $uH$ is the
single element $u$, so $G/H=SU(2)$, and the correspondence of (17.10) is
\begin{equation}\label{quat:first-column}
 F\colon SU(2)\to S^3,\qquad F(u)=u\begin{bmatrix}1\\0\end{bmatrix}
 =\begin{bmatrix}u_{11}\\u_{21}\end{bmatrix}.
\end{equation}
Its inverse is the matrix of Eq.~\eqref{quat:frankel-19-11},
$F^{-1}(z_1,z_2)=\begin{bmatrix}z_1&-\bar z_2\\z_2&\bar z_1\end{bmatrix}$.
Check both compositions. $F(F^{-1}(z_1,z_2))$ is the first column
$(z_1,z_2)^T$. For $u$ in the form~\eqref{quat:su2-general},
$F(u)=(\alpha,\beta)^T$ and $F^{-1}(\alpha,\beta)=u$. In real coordinates the
entries of $F$ and $F^{-1}$ are linear in the real and imaginary parts of
their arguments, so both maps are continuous, indeed smooth. Therefore $F$ is
a homeomorphism:
\begin{equation}
 SU(2)\approx S^3 .
\end{equation}
The dimensions agree: $S^3$ is cut out of $\mathbb R^4$ by one equation,
giving $4-1=3=\dim SU(2)$.

In quaternion form, put $\alpha=q_0-iq_3$ and $\beta=q_2-iq_1$. By
Eq.~\eqref{quat:minus-i-sigma},
\begin{equation}
 q_0I+q^a(-i\sigma_a)
 =\begin{bmatrix}q_0-iq_3&-iq_1-q_2\\-iq_1+q_2&q_0+iq_3\end{bmatrix}
 =\begin{bmatrix}\alpha&-\bar\beta\\ \beta&\bar\alpha\end{bmatrix},
\end{equation}
and $|\alpha|^2+|\beta|^2=q_0^2+q_1^2+q_2^2+q_3^2=|q|^2$. Under
Eq.~\eqref{quat:dictionary}, $SU(2)$ is the group of unit quaternions, the
unit sphere $S^3\subset\mathbb H\cong\mathbb R^4$. Holm writes the same
matrix as $Q$ (Eq.~(3.2.11), p.~100).

\subsection{Connectedness, and \texorpdfstring{$\det\Phi(u)=+1$}{det Phi(u)=+1}}
A topological space $X$ is \emph{connected} if it is not the union of two
disjoint nonempty open subsets. Equivalently, every continuous map
$f\colon X\to\{-1,+1\}$ is constant. If $f$ took both values, then
$f^{-1}(-1)$ and $f^{-1}(+1)$ would split $X$; if $X=U\cup V$ were split,
$f=-1$ on $U$ and $f=+1$ on $V$ would be continuous and not constant.

$X$ is \emph{path-connected} if any two points are joined by a continuous
path $\gamma\colon[0,1]\to X$. Path-connected implies connected: if
$f\colon X\to\{\pm1\}$ is continuous, then $f\circ\gamma\colon[0,1]\to
\{\pm1\}$ is continuous, and by the intermediate value theorem it cannot
pass from $-1$ to $+1$, so $f(\gamma(0))=f(\gamma(1))$.

$S^3$ is path-connected. Let $p,q\in S^3$ with $q\ne-p$. If
$(1-t)p+tq=0$ for some $t\in[0,1]$, then taking lengths gives $1-t=t$, so
$t=\frac12$ and $q=-p$. Hence the segment from $p$ to $q$ avoids $0$, and
\begin{equation}
 \gamma(t)=\frac{(1-t)p+tq}{|(1-t)p+tq|}
\end{equation}
is a continuous path in $S^3$ from $p$ to $q$. If $q=-p$, choose
$r\in S^3$ with $r\ne\pm p$ and join $p$ to $r$ and then $r$ to $-p$. Thus
$S^3$ is connected, and so is $SU(2)$, its image under the continuous map
$F^{-1}$. This is Frankel's ``in particular $SU(2)=S^3$ is connected''.

Now apply this to $\det\Phi(u)$. By Eq.~\eqref{quat:Phi-entries}, the entries
of $\Phi(u)$ are polynomials in the real and imaginary parts of the entries
of $u$, using $u^{-1}=u^\dagger$. So $u\mapsto\det\Phi(u)$ is continuous on
$SU(2)$, with values in $\{\pm1\}$ because $\Phi(u)$ is orthogonal. On the
connected $SU(2)$ it is constant, and at $u=I$ we have $\Phi(I)=I$, with
determinant $1$. Therefore
\begin{equation}
 \Phi\colon SU(2)\to SO(3).
\end{equation}
By contrast, $O(3)$ is not connected: $\det\colon O(3)\to\{\pm1\}$ is
continuous and takes both values, at $I$ and at $-I$.

\section{\texorpdfstring{$\Phi\colon SU(2)\to SO(3)$}{Phi: SU(2) to SO(3)} is onto and two-to-one}\label{quat:sec:covering}
This section follows Frankel's Section~19.1d (pp.~496--497).

\begin{theorem}[Frankel (19.12)]\label{quat:thm:covering}
$\Phi\colon SU(2)\to SO(3)$ is onto. For each $R\in SO(3)$ there are
exactly two elements $\pm u\in SU(2)$ with $\Phi(\pm u)=R$.
\end{theorem}

Frankel states this for $\mathrm{Ad}$ of his Eq.~(19.10), which is our
$\Phi$. In his proof $h$ denotes both a vector of $\mathbb R^3$ and the
Hermitian matrix $h^j\sigma_j$. We write $\mathbf h\in\mathbb R^3$ and
$\Sigma(\mathbf h)$. His $\mathrm{Ad}_*$ is the derivative of $\Phi$ at $I$,
which we write $d\Phi$.

\subsection{The derivative of \texorpdfstring{$\Phi$}{Phi}}
\emph{One-parameter subgroups of $SU(2)$.} The argument of
Eqs.~\eqref{quat:one-param-ode}--\eqref{quat:one-param-exp} works verbatim in
$GL(2,\mathbb C)$, so every differentiable one-parameter subgroup of
$SU(2)$ is $u(t)=\exp(tX)$ with $X\in\mathfrak{su}(2)$. By
Eq.~\eqref{quat:su2-bases}, $X=-i\Sigma(\mathbf h)$ for a unique
$\mathbf h\in\mathbb R^3$:
\begin{equation}\label{quat:u-of-t}
 u(t)=\exp\bigl(-it\,\Sigma(\mathbf h)\bigr)
 =\exp\Bigl(\frac{t}{i}\,h^j\sigma_j\Bigr).
\end{equation}
This is Frankel's $u(t)=e^{th/i}=e^{-ith}$.

\emph{Its image is a one-parameter subgroup of $SO(3)$.} Since $\Phi$ is a
homomorphism, $\Phi(u(t+s))=\Phi(u(t)u(s))=\Phi(u(t))\Phi(u(s))$ and
$\Phi(u(0))=I$. By Eq.~\eqref{quat:one-param-exp},
\begin{equation}
 \Phi\bigl(u(t)\bigr)=\exp(tS),\qquad
 S=\frac{d}{dt}\Phi\bigl(u(t)\bigr)\Big|_{t=0}=:d\Phi(X).
\end{equation}

\emph{Computing $S$.} Apply Eq.~\eqref{quat:Phi} with $u=u(t)$:
\begin{equation}
 \Sigma\bigl(\Phi(u(t))\mathbf x\bigr)
 =e^{-it\Sigma(\mathbf h)}\,\Sigma(\mathbf x)\,e^{it\Sigma(\mathbf h)} .
\end{equation}
Differentiate at $t=0$. On the left, $\Sigma$ is linear and does not depend
on $t$, so the derivative is $\Sigma(S\mathbf x)$. On the right, the product
rule with $\frac{d}{dt}e^{\mp it\Sigma(\mathbf h)}\big|_{t=0}
=\mp i\Sigma(\mathbf h)$ gives
\begin{equation}
 \Sigma(S\mathbf x)=-i\Sigma(\mathbf h)\Sigma(\mathbf x)
   +\Sigma(\mathbf x)\,i\Sigma(\mathbf h)
 =-i\bigl[\Sigma(\mathbf h),\Sigma(\mathbf x)\bigr].
\end{equation}
Expand the commutator with Frankel's Eq.~(19.6),
$[\sigma_j,\sigma_k]=2i\epsilon_{jkl}\sigma_l$, and
$(\mathbf h\times\mathbf x)^l=\epsilon_{ljk}h^jx^k=\epsilon_{jkl}h^jx^k$:
\begin{align}
 -i\bigl[\Sigma(\mathbf h),\Sigma(\mathbf x)\bigr]
 &=-ih^jx^k[\sigma_j,\sigma_k]
  =-ih^jx^k\,2i\,\epsilon_{jkl}\sigma_l\notag\\
 &=2\epsilon_{jkl}h^jx^k\sigma_l
  =2(\mathbf h\times\mathbf x)^l\sigma_l
  =\Sigma(2\mathbf h\times\mathbf x).\label{quat:dPhi-computation}
\end{align}
So $\Sigma(S\mathbf x)=\Sigma(2\mathbf h\times\mathbf x)$ for every
$\mathbf x$. Since $\Sigma$ is injective, $S\mathbf x=2\mathbf h\times\mathbf x$,
that is $S=\widehat{2\mathbf h}$ by Eq.~\eqref{quat:hat}. The
one-parameter group $\Phi(u(t))$ rotates with angular velocity
$\boldsymbol\omega=2\mathbf h$, and Eq.~\eqref{quat:frankel-19-3} gives
\begin{equation}\label{quat:frankel-19-13}
 \Phi\Bigl(\exp\bigl(-it\,\Sigma(\mathbf h)\bigr)\Bigr)
 =\exp\bigl(t\,(2h)^jE_j\bigr).
\end{equation}
This is Frankel's Eq.~(19.13),
$\mathrm{Ad}\exp(\frac{\boldsymbol\sigma}{i}\bullet\mathbf ht)\mathbf x_*
\sim R(t)\mathbf x=\exp(\mathbf E\bullet2\mathbf ht)\mathbf x$. Each $\sim$
there is an instance of Eq.~\eqref{quat:Phi}: the matrix on the left is
$\Sigma$ of the vector on the right.

\emph{$d\Phi$ is $\mathrm{ad}$.} The first equality of
Eq.~\eqref{quat:dPhi-computation} reads, with $X=-i\Sigma(\mathbf h)$,
\begin{equation}
 \Sigma\bigl(d\Phi(X)\mathbf x\bigr)=\bigl[X,\Sigma(\mathbf x)\bigr]
 =\mathrm{ad}_X\Sigma(\mathbf x).
\end{equation}
So $d\Phi(X)$ is $\mathrm{ad}_X$ written in the coordinates $\Sigma$, just as
$\Phi(u)$ is $\mathrm{Ad}_u$ written in those coordinates. The identity
$\Phi(\exp tX)=\exp(t\,d\Phi(X))$ is the relation
$\mathrm{Ad}_{\exp(tX)}=\exp(t\,\mathrm{ad}_X)$ of
\S\ref{dgdt:sec:adjoint-representation}, read in $\mathbb R^3$.

\emph{Frankel's Eq.~(19.14).} Take $X=\sigma_a/(2i)=-i\Sigma(\frac12\mathbf e_a)$.
Then $\mathbf h=\frac12\mathbf e_a$, $2\mathbf h=\mathbf e_a$, and
$\widehat{\mathbf e_a}=E_a$:
\begin{equation}\label{quat:frankel-19-14}
 d\Phi\Bigl(\frac{\sigma_a}{2i}\Bigr)=E_a .
\end{equation}
Frankel calls this unsurprising given the remarks after his Eq.~(19.6): the
basis $\sigma_a/(2i)$ has the same structure constants as $E_a$
(Eq.~\eqref{quat:Fa-brackets}). For Kosmann-Schwarzbach's basis,
$d\Phi(\xi_1)=-E_1$, $d\Phi(\xi_2)=E_2$, $d\Phi(\xi_3)=-E_3$ by
Eq.~\eqref{quat:KS-vs-Frankel}. In either case $d\Phi$ carries a basis to a
basis, so $d\Phi\colon\mathfrak{su}(2)\to\mathfrak{so}(3)$ is a linear
isomorphism between spaces of dimension $3$.

\subsection{Example: rotations about the \texorpdfstring{$z$}{z} axis}
Take $\mathbf h=\mathbf e_3$, so $\Sigma(\mathbf h)=\sigma_3$; Frankel
writes $\mathbf h=\boldsymbol\sigma(0,0,1)^T$, that is $\mathbf h=\sigma_3$.
Since $\sigma_3/i=-i\sigma_3=\operatorname{diag}(-i,i)$ is diagonal, every
power is diagonal and the exponential is taken entry by entry:
\begin{equation}
 u(\theta)=\exp\Bigl(\theta\frac{\sigma_3}{i}\Bigr)
 =\begin{bmatrix}e^{-i\theta}&0\\0&e^{i\theta}\end{bmatrix}.
\end{equation}
By Eq.~\eqref{quat:frankel-19-13} with $t=\theta$,
$\Phi(u(\theta))=\exp(2\theta E_3)$, which is $R(2\theta)$ by
Eq.~\eqref{quat:exp-E3} with $\theta$ replaced by $2\theta$.

The same result follows directly from Eq.~\eqref{quat:Phi}, without
derivatives. Multiplying on the left by $u(\theta)$ multiplies row~1 by
$e^{-i\theta}$ and row~2 by $e^{i\theta}$. Multiplying on the right by
$u(\theta)^{-1}=\operatorname{diag}(e^{i\theta},e^{-i\theta})$ multiplies
column~1 by $e^{i\theta}$ and column~2 by $e^{-i\theta}$:
\begin{equation}
 u(\theta)\Sigma(\mathbf x)u(\theta)^{-1}
 =\begin{bmatrix}z&e^{-2i\theta}(x-iy)\\e^{2i\theta}(x+iy)&-z\end{bmatrix}.
\end{equation}
This is $\Sigma(\mathbf x')$ with $z'=z$ and
$x'+iy'=e^{2i\theta}(x+iy)=(x\cos2\theta-y\sin2\theta)
+i(x\sin2\theta+y\cos2\theta)$. So $\Phi(u(\theta))=R(2\theta)$.

As $\theta$ runs from $0$ to $2\pi$, $e^{-i\theta}$ goes once around the unit
circle, so $u(\theta)$ is a simple closed curve in $SU(2)$, passing through
$u(\pi)=-I$. Its image $R(2\theta)$ makes two full rotations, returning to
$I$ at $\theta=\pi$ and at $\theta=2\pi$. The half $0\le\theta\le\pi$ is a
path from $I$ to $-I$ in $SU(2)$ whose image is a closed loop in $SO(3)$; that
loop is the one studied in Frankel's Section~19.2a and in his
Section~21.2.

\subsection{\texorpdfstring{$\Phi$}{Phi} is onto}
Frankel uses that every rotation is a rotation about some axis. Let
$R\in SO(3)$. Then
\begin{align}
 \det(R-I)&=\det\bigl(R(I-R^T)\bigr)=\det R\,\det\bigl((I-R)^T\bigr)\notag\\
 &=\det(I-R)=(-1)^3\det(R-I),
\end{align}
using $RR^T=I$, $\det R=1$, and $\det A^T=\det A$. So $\det(R-I)=0$, and $R$
has a unit eigenvector $\mathbf n$ with $R\mathbf n=\mathbf n$. $R$ preserves
the plane perpendicular to $\mathbf n$: if $\mathbf x\cdot\mathbf n=0$ then
$R\mathbf x\cdot\mathbf n=R\mathbf x\cdot R\mathbf n=\mathbf x\cdot\mathbf n=0$.
In a right-handed orthonormal basis $(\mathbf u,\mathbf n\times\mathbf u,
\mathbf n)$, the matrix of $R$ is block diagonal with an upper $2\times2$
block $A$ and a $1$ in the corner. $A$ has orthonormal columns and
$\det A=\det R=1$. Its first column is a unit vector
$(\cos\theta,\sin\theta)$ for some $\theta$. The second is a unit vector
orthogonal to it, $\pm(-\sin\theta,\cos\theta)$, and $\det A=1$ selects the
plus sign. So $R=R_{\mathbf n}(\theta)=\exp(\theta\,n^jE_j)$ by
Eq.~\eqref{quat:frankel-19-4}.

Now take $\mathbf h=\frac12\mathbf n$ and $t=\theta$ in
Eq.~\eqref{quat:frankel-19-13}:
\begin{equation}\label{quat:onto}
 \Phi\Bigl(\exp\Bigl(-i\frac{\theta}{2}\Sigma(\mathbf n)\Bigr)\Bigr)
 =\exp\bigl(\theta\,n^jE_j\bigr)=R .
\end{equation}
This is Frankel's $\mathrm{Ad}\exp(\boldsymbol\sigma/2i\bullet\boldsymbol\omega
\theta)=R$, with his $\boldsymbol\omega$ the unit axis $\mathbf n$. By
Eq.~\eqref{quat:half-angle}, the element of $SU(2)$ in Eq.~\eqref{quat:onto}
is $\cos\frac{\theta}{2}I-i\sin\frac{\theta}{2}\Sigma(\mathbf n)$, the unit
quaternion of Eq.~\eqref{quat:axis-angle}. So $\Phi$ is onto.

\subsection{\texorpdfstring{$\Phi$}{Phi} in closed form: onto and two-to-one explicitly}\label{quat:sec:Phi-closed}
Every $u\in SU(2)$ is $u=q_0I-i\Sigma(\mathbf q)$ with
$q_0^2+|\mathbf q|^2=1$ (the quaternion form in
\S\ref{quat:sec:su2-sphere}, since $q^a(-i\sigma_a)=-i\Sigma(\mathbf q)$), and
$u^{-1}=u^\dagger=q_0I+i\Sigma(\mathbf q)$ because $\Sigma(\mathbf q)$ is
Hermitian. Computing $\Phi(u)$ from Eq.~\eqref{quat:Phi} needs one product
rule for Pauli matrices. Adding half of Frankel's anticommutator
(Eq.~(19.18)) and half of his commutator (Eq.~(19.6)) gives
$\sigma_j\sigma_k=\delta_{jk}I+i\epsilon_{jkl}\sigma_l$, so
\begin{equation}\label{quat:Sigma-product}
 \Sigma(\mathbf a)\Sigma(\mathbf b)=a^jb^k\sigma_j\sigma_k
 =(\mathbf a\cdot\mathbf b)\,I+i\,\Sigma(\mathbf a\times\mathbf b),
\end{equation}
Frankel's Eq.~(19.19), with $(\mathbf a\times\mathbf b)^l=\epsilon_{jkl}a^jb^k$.
Its antisymmetric part is
$[\Sigma(\mathbf a),\Sigma(\mathbf b)]=2i\,\Sigma(\mathbf a\times\mathbf b)$.
Now expand
\begin{equation}
 u\Sigma(\mathbf x)u^{-1}
 =q_0^2\Sigma(\mathbf x)
  +iq_0\bigl[\Sigma(\mathbf x),\Sigma(\mathbf q)\bigr]
  +\Sigma(\mathbf q)\Sigma(\mathbf x)\Sigma(\mathbf q),
\end{equation}
and treat the last two terms separately. The middle term is
$iq_0\cdot2i\,\Sigma(\mathbf x\times\mathbf q)=2q_0\Sigma(\mathbf q\times\mathbf x)$.
For the last term apply Eq.~\eqref{quat:Sigma-product} twice. The scalar
$(\mathbf q\times\mathbf x)\cdot\mathbf q$ vanishes, and
$(\mathbf q\times\mathbf x)\times\mathbf q=|\mathbf q|^2\mathbf x
-(\mathbf q\cdot\mathbf x)\mathbf q$ by
$(\mathbf a\times\mathbf b)\times\mathbf c=(\mathbf a\cdot\mathbf c)\mathbf b
-(\mathbf b\cdot\mathbf c)\mathbf a$:
\begin{align}
 \Sigma(\mathbf q)\Sigma(\mathbf x)\Sigma(\mathbf q)
 &=\bigl((\mathbf q\cdot\mathbf x)I+i\Sigma(\mathbf q\times\mathbf x)\bigr)
   \Sigma(\mathbf q)\notag\\
 &=(\mathbf q\cdot\mathbf x)\Sigma(\mathbf q)
   -\Sigma\bigl((\mathbf q\times\mathbf x)\times\mathbf q\bigr)\notag\\
 &=2(\mathbf q\cdot\mathbf x)\Sigma(\mathbf q)-|\mathbf q|^2\Sigma(\mathbf x).
\end{align}
Collecting terms inside $\Sigma$ and using Eq.~\eqref{quat:Phi},
\begin{equation}\label{quat:euler-rodrigues}
 \Phi(u)\mathbf x=\bigl(q_0^2-|\mathbf q|^2\bigr)\mathbf x
 +2q_0\,\mathbf q\times\mathbf x+2(\mathbf q\cdot\mathbf x)\,\mathbf q .
\end{equation}
In matrix form this is
$\Phi(u)=(q_0^2-|\mathbf q|^2)I+2q_0\widehat{\mathbf q}+2\mathbf q\mathbf q^T$,
the Euler--Rodrigues formula.

\emph{Two-to-one, explicitly.} Every term of Eq.~\eqref{quat:euler-rodrigues}
is quadratic in $(q_0,\mathbf q)$. Replacing $u$ by $-u$ replaces
$(q_0,\mathbf q)$ by $(-q_0,-\mathbf q)$ and leaves $\Phi(u)$ unchanged.
The kernel can also be read off. If $\Phi(u)=I$ and $\mathbf q\ne0$, take a
unit $\mathbf x\perp\mathbf q$. Then Eq.~\eqref{quat:euler-rodrigues} gives
$(q_0^2-|\mathbf q|^2)\mathbf x+2q_0\,\mathbf q\times\mathbf x=\mathbf x$. The
two terms are perpendicular, so $q_0^2-|\mathbf q|^2=1$, which together with
$q_0^2+|\mathbf q|^2=1$ forces $\mathbf q=0$, a contradiction. So
$\mathbf q=0$, $q_0=\pm1$, and $u=\pm I$.

\emph{Onto, explicitly.} Given $R=R_{\mathbf n}(\theta)$, which covers every
rotation by the axis argument above, put $q_0=\cos\frac{\theta}{2}$ and
$\mathbf q=\sin\frac{\theta}{2}\mathbf n$. The double-angle formulas give
$q_0^2-|\mathbf q|^2=\cos\theta$, $2q_0\mathbf q=\sin\theta\,\mathbf n$, and
$2\mathbf q\mathbf q^T=(1-\cos\theta)\mathbf n\mathbf n^T$. Then
Eq.~\eqref{quat:euler-rodrigues} becomes
\begin{equation}
 \Phi(u)\mathbf x=\cos\theta\,\mathbf x+\sin\theta\,\mathbf n\times\mathbf x
 +(1-\cos\theta)(\mathbf n\cdot\mathbf x)\mathbf n ,
\end{equation}
which is Eq.~\eqref{quat:rotate-about-n} after writing
$\mathbf x=(\mathbf n\cdot\mathbf x)\mathbf n+\mathbf x_\perp$. So the two
preimages of the rotation by $\theta$ about $\mathbf n$ are
\begin{equation}\label{quat:two-preimages}
 \pm u=\pm\Bigl(\cos\frac{\theta}{2}\,I-i\sin\frac{\theta}{2}\,\Sigma(\mathbf n)\Bigr).
\end{equation}
The minus sign is the same formula with $\theta+2\pi$ in place of $\theta$,
since $\cos\frac{\theta+2\pi}{2}=-\cos\frac{\theta}{2}$ and
$\sin\frac{\theta+2\pi}{2}=-\sin\frac{\theta}{2}$. The rotation angles
$\theta$ and $\theta+2\pi$ give the same rotation of $\mathbb R^3$ but
different elements of $SU(2)$.

\subsection{\texorpdfstring{$\Phi$}{Phi} is two-to-one, and \texorpdfstring{$SO(3)\approx\mathbb RP^3$}{SO(3) = RP3}}
For every $u$, $(-u)\Sigma(\mathbf x)(-u)^{-1}
=(-1)(-1)^{-1}u\Sigma(\mathbf x)u^{-1}=u\Sigma(\mathbf x)u^{-1}$, so
$\Phi(-u)=\Phi(u)$. $\Phi$ is therefore at least two-to-one; in Frankel's
words, the representation is not faithful.

\emph{The kernel.} $\Phi(u)=I$ means $u\Sigma(\mathbf x)u^{-1}=\Sigma(\mathbf x)$
for all $\mathbf x$. Taking $\mathbf x=\mathbf e_k$, and conversely by
linearity, this holds exactly when $u$ commutes with $\sigma_1,\sigma_2,\sigma_3$.
Since $\sigma_k^{-1}=\sigma_k$, commuting with $\sigma_k$ is
$\sigma_ku\sigma_k=u$, Frankel's $\sigma_\alpha^{-1}u\sigma_\alpha=u$. Write $u$
in the form~\eqref{quat:su2-general}. Conjugating by
$\sigma_3=\operatorname{diag}(1,-1)$ changes the sign of the off-diagonal
entries, and conjugating by $\sigma_1$ exchanges both the rows and the
columns:
\begin{equation}
 \sigma_3u\sigma_3=\begin{bmatrix}\alpha&\bar\beta\\-\beta&\bar\alpha\end{bmatrix},
 \qquad
 \sigma_1u\sigma_1=\begin{bmatrix}\bar\alpha&\beta\\-\bar\beta&\alpha\end{bmatrix}.
\end{equation}
The first equals $u$ only if $-\bar\beta=\bar\beta$, so $\beta=0$. The
second equals $u$ only if $\bar\alpha=\alpha$, so $\alpha$ is real. Then
$|\alpha|^2+|\beta|^2=1$ gives $\alpha=\pm1$, and
\begin{equation}
 \ker\Phi=\{I,-I\}.
\end{equation}

\emph{The fibers.} $\Phi(u)=\Phi(u')$ if and only if
$\Phi(u^{-1}u')=\Phi(u)^{-1}\Phi(u')=I$, if and only if
$u^{-1}u'\in\{\pm I\}$, if and only if $u'=\pm u$. Since $u\ne-u$, every
rotation has exactly two preimages. This completes the proof of
Theorem~\ref{quat:thm:covering}.

\emph{The elementary result of group theory.} Frankel continues: ``It is an
elementary result of group theory that if $\phi\colon G\to G'$ is a
homomorphism of $G$ onto $G'$, then $G'$ is isomorphic to the coset space
$G/H$, where $H=\phi^{-1}(e')$ is the kernel'' (pp.~496--497). This is the first
isomorphism theorem for groups. The algebra part proves the same theorem for
modules (Theorem~\ref{agat:Thm:FirstIsoModules}, from Rotman). Modules are
abelian groups, where every subgroup is normal; for groups, normality of the
kernel is the extra step. Dummit and Foote state the group version as
Theorem~16 of their Section~3.3 \cite{agat:DuFo2003}. Quotient groups appear
in Jordan's and H\"older's work on composition series (1870--1889); the
isomorphism theorems in their general form are usually credited to Emmy
Noether (1927), and they reached textbooks through van der Waerden's
\emph{Moderne Algebra} (1930).

\begin{theorem}[first isomorphism theorem]\label{quat:thm:first-iso}
Let $\varphi\colon G\to G'$ be a group homomorphism and
$K=\varphi^{-1}(e')$ its kernel. Then $K$ is a normal subgroup of $G$, the
cosets $G/K$ form a group under $(gK)(hK)=ghK$, and
\begin{equation}
 \bar\varphi\colon G/K\to\varphi(G),\qquad \bar\varphi(gK)=\varphi(g),
\end{equation}
is a well-defined group isomorphism. If $\varphi$ is onto, $G/K\cong G'$.
\end{theorem}
\begin{proof}
\emph{$K$ is normal.} For $k\in K$ and $g\in G$,
$\varphi(gkg^{-1})=\varphi(g)\,e'\,\varphi(g)^{-1}=e'$, so $gKg^{-1}\subseteq K$.

\emph{The coset product is well defined.} If $g'=gk$ and $h'=hk'$ with
$k,k'\in K$, then $g'h'=gh\,(h^{-1}kh)\,k'$, and $h^{-1}kh\in K$ by
normality, so $g'h'K=ghK$.

\emph{$\bar\varphi$ is well defined and injective.} For $g,h\in G$,
\begin{equation}
 gK=hK\iff g^{-1}h\in K\iff\varphi(g^{-1}h)=e'\iff\varphi(g)=\varphi(h).
\end{equation}
Reading the chain left to right shows that $\bar\varphi(gK)$ does not depend
on the representative $g$; reading it right to left shows that $\bar\varphi$
is injective.

\emph{$\bar\varphi$ is a homomorphism onto $\varphi(G)$.}
$\bar\varphi(gK\,hK)=\bar\varphi(ghK)=\varphi(gh)=\varphi(g)\varphi(h)$, and
every $\varphi(g)$ equals $\bar\varphi(gK)$.
\end{proof}
Frankel derives the same correspondence from his Fundamental
Principle~(17.10). $G$ acts on $G'$ by $g\cdot g'=\varphi(g)g'$. The action
is transitive when $\varphi$ is onto: to move $g_1'$ to $g_2'$, choose $g$
with $\varphi(g)=g_2'g_1'^{-1}$. The stability subgroup of $e'$ is
$\ker\varphi$. That argument gives the bijection
$G/\ker\varphi\leftrightarrow G'$ of sets.
Theorem~\ref{quat:thm:first-iso} adds that the bijection respects
multiplication. For $\Phi$, which is onto with kernel $\{\pm I\}$,
\begin{equation}\label{quat:SO3-quotient}
 SU(2)/\{\pm I\}\cong SO(3)\quad\text{as groups},\qquad
 \bar\Phi\bigl(\{u,-u\}\bigr)=\Phi(u).
\end{equation}

In the coordinates of Eq.~\eqref{quat:first-column},
$F(-u)=-F(u)$, so each fiber $\{u,-u\}$ is a pair of antipodal points of
$S^3$. Identifying antipodal points of $S^3$ gives the real projective space
$\mathbb RP^3=S^3/\{\pm1\}$, the space of lines through $0$ in $\mathbb R^4$,
since each such line meets $S^3$ in one antipodal pair:
\begin{equation}
\begin{tikzcd}
 S^3 \arrow[r,"\approx"] \arrow[d,"p"'] & SU(2) \arrow[d,"\Phi"]\\
 \mathbb RP^3 \arrow[r,"\approx"] & SO(3)
\end{tikzcd}
\end{equation}
Here $p(x)=\{x,-x\}$. The induced map $\mathbb RP^3\to SO(3)$ is a
continuous bijection. $\mathbb RP^3$ is compact, as the image of the compact
$S^3$ under $p$, and $SO(3)\subset\mathbb R^9$ is Hausdorff, so the map is a
homeomorphism: $SO(3)\approx\mathbb RP^3$. By Theorem~(17.11), $SU(2)$ is a
principal bundle over $SU(2)/\{\pm I\}=SO(3)$ with the two-point fiber
$\{\pm I\}$, that is, a two-sheeted covering. Its dimension count is
$\dim SO(3)=\dim SU(2)-\dim\{\pm I\}=3-0$. Also, $SO(3)=\Phi(SU(2))$ is
connected, as the continuous image of a connected space.

\section{The smooth statements, with Lee's theorems}\label{quat:sec:smooth-cover}
The preceding sections prove statements about sets, groups and continuous
maps. Frankel's ``$\approx$'' and ``principal bundle'' need the smooth
versions. Lee's \emph{Introduction to Smooth Manifolds}, second edition
\cite{dgdt:JLee2012}, supplies each one; below, every step names the result
it uses, with its printed page.

\subsection{\texorpdfstring{$S^3$}{S3} and \texorpdfstring{$SU(2)$}{SU(2)} are smooth manifolds}
Let $f(x)=|x|^2$ on $\mathbb R^4$. Then $df_x=2x^T\ne0$ on $S^3=f^{-1}(1)$,
so $1$ is a regular value. By the regular level set theorem
(Lee, Corollary~5.14, p.~106), $S^3$ is a properly embedded submanifold of
codimension $1$, of dimension $3$. $SU(2)$ is a properly embedded
$(2^2-1)$-dimensional Lie subgroup of $U(2)$, hence embedded in
$GL(2,\mathbb C)$ (Lee, Example~7.30, p.~168, and Problem~7-14, p.~172).
Likewise $SO(3)$ is an embedded Lie subgroup of $GL(3,\mathbb R)$ (Lee,
Example~7.28, p.~167).

Two facts about restrictions are used repeatedly. If $G\colon M\to N$ is
smooth and $S\subseteq M$ is an embedded submanifold, then $G|_S$ is smooth
(Lee, Theorem~5.27, p.~112). If $G\colon N\to M$ is smooth with image in an
embedded submanifold $S\subseteq M$, then $G$ is smooth as a map into $S$
(Lee, Corollary~5.30, p.~113).

\subsection{\texorpdfstring{$SU(2)$}{SU(2)} is diffeomorphic to \texorpdfstring{$S^3$}{S3}}
This is Lee's Problem~7-16 (p.~172). The map $F$ of
Eq.~\eqref{quat:first-column} is the restriction to $SU(2)$ of the linear map
$M(2,\mathbb C)\to\mathbb C^2$, $u\mapsto u(1,0)^T$, with image in $S^3$. Its
inverse is the restriction to $S^3$ of the real-linear map
$\mathbb C^2\to M(2,\mathbb C)$ of Eq.~\eqref{quat:frankel-19-11}, with image
in $SU(2)$. Linear maps are smooth. By the two restriction facts, $F$ and
$F^{-1}$ are smooth, so $F$ is a diffeomorphism:
\begin{equation}\label{quat:SU2-diffeo-S3}
 SU(2)\approx S^3\quad\text{(diffeomorphic)}.
\end{equation}

\subsection{\texorpdfstring{$\Phi$}{Phi} is a smooth double covering}
By Eq.~\eqref{quat:Phi-entries}, with $u^{-1}=u^\dagger$, each entry of
$\Phi(u)$ is a polynomial in the real and imaginary parts of the entries of
$u$. So $\Phi$ is the restriction of a smooth map $M(2,\mathbb C)\to
M(3,\mathbb R)$, with image in $SO(3)$, and the restriction facts make
$\Phi\colon SU(2)\to SO(3)$ smooth. It is a homomorphism, hence a Lie group
homomorphism. Every Lie group homomorphism has constant rank (Lee,
Theorem~7.5, p.~153). By Eq.~\eqref{quat:frankel-19-14}, $d\Phi$ at $I$ is an
isomorphism, so $d\Phi_u$ has rank $3$ at every $u$.

Lee's Theorem~7.5 is a short computation. With $L_u$ left multiplication,
$\Phi\circ L_u=L_{\Phi(u)}\circ\Phi$ because $\Phi$ is a homomorphism. The
chain rule gives $d\Phi_u\circ d(L_u)_I=d(L_{\Phi(u)})_I\circ d\Phi_I$, and
both $d(L_u)_I$ and $d(L_{\Phi(u)})_I$ are isomorphisms, since $L_u$ and
$L_{\Phi(u)}$ are diffeomorphisms. So $d\Phi_u$ and $d\Phi_I$ have the same
rank.

Both groups are connected (\S\ref{quat:sec:su2-sphere} and
\S\ref{quat:sec:covering}). Lee's Theorem~21.31 (Homomorphisms with Discrete
Kernels, p.~557) states that for connected Lie groups $G,H$ and a Lie group
homomorphism $F\colon G\to H$ the following are equivalent: (a) $F$ is
surjective and has discrete kernel; (b) $F$ is a smooth covering map; (c) $F$
is a local diffeomorphism; (d) $F_*\colon\operatorname{Lie}(G)\to
\operatorname{Lie}(H)$ is an isomorphism. Here (a) is
Theorem~\ref{quat:thm:covering} with $\ker\Phi=\{\pm I\}$, and (d) is
Eq.~\eqref{quat:frankel-19-14}; either one gives (b):
\begin{equation}\label{quat:Phi-smooth-cover}
 \Phi\colon SU(2)\to SO(3)\ \text{is a smooth two-sheeted covering map.}
\end{equation}
This is the precise content of Frankel's ``principal bundle with two-point
fiber''.

\subsection{\texorpdfstring{$\mathbb RP^3\approx SU(2)/\{\pm I\}\cong SO(3)$}{RP3 = SU(2)/{+-I} = SO(3)}}
Theorem~\ref{quat:thm:first-iso} identifies $SU(2)/\{\pm I\}$ with $SO(3)$ as
groups. Lee's First Isomorphism Theorem for Lie Groups (Theorem~21.27,
p.~556) adds the smooth structure: the kernel $\{\pm I\}$ is a closed normal
Lie subgroup, and $\Phi$ descends to a Lie group isomorphism
\begin{equation}
 \bar\Phi\colon SU(2)/\{\pm I\}\to SO(3),
\end{equation}
that is, a group isomorphism that is also a diffeomorphism. The quotient
$SU(2)/\{\pm I\}$ carries the smooth structure of Lee's Theorem~21.29
(p.~556): $\{\pm I\}$ is a discrete subgroup of the connected $SU(2)$, so the
quotient map is a smooth normal covering map.

For $\mathbb RP^3$, the group $\{\pm1\}$ acts on $S^3$ by $x\mapsto\pm x$. The
action is smooth and free, and it is proper because $\{\pm1\}$ is compact
(Lee, Corollary~21.6, p.~544). Lee's Theorem~21.13 (p.~549) makes
$S^3/\{\pm1\}=\mathbb RP^3$ a smooth manifold with $p\colon S^3\to\mathbb RP^3$
a smooth normal covering map; for $\mathbb RP^n$ with its standard structure
this is Lee's Problem~4-10 (p.~96). Since $F(-u)=-F(u)$, the composite
$S^3\xrightarrow{F^{-1}}SU(2)\to SU(2)/\{\pm I\}$ is constant on the fibers
of $p$. A smooth covering map is a surjective smooth submersion, so the
composite descends to a smooth map $\mathbb RP^3\to SU(2)/\{\pm I\}$ (Lee,
Theorem~4.30, Passing Smoothly to the Quotient, p.~90). The same argument
with $F$ in place of $F^{-1}$ gives its smooth inverse. Composing with
$\bar\Phi$,
\begin{equation}\label{quat:SO3-RP3-smooth}
 \mathbb RP^3\;\approx\;SU(2)/\{\pm I\}\;\cong\;SO(3),
\end{equation}
a diffeomorphism. This upgrades the homeomorphism of the previous section.

This is Lee's Problem~21-21 (p.~562), in quaternion language. There $S$ is
the group of unit quaternions, $E$ the imaginary quaternions, and $\rho(q)$
the matrix of $v\mapsto qvq^*$ on $E$ in the basis $(\mathbf i,\mathbf
j,\mathbf k)$. Under the dictionary~\eqref{quat:dictionary},
$e_a\leftrightarrow-i\sigma_a$, the imaginary quaternion $x^ae_a$ corresponds
to $-i\Sigma(\mathbf x)$, and $q^*=q^{-1}$ for a unit quaternion. Multiplying
Eq.~\eqref{quat:Phi} by $-i$ gives
$u\bigl(-i\Sigma(\mathbf x)\bigr)u^{-1}=-i\Sigma\bigl(\Phi(u)\mathbf x\bigr)$.
So $\Phi(u)$ is Lee's $\rho(q)$, and the problem's parts (a)--(c) are
Eqs.~\eqref{quat:Phi}, \eqref{quat:euler-rodrigues} and
\eqref{quat:SO3-RP3-smooth}.

\subsection{\texorpdfstring{$\pi_1(SO(3))=\mathbb Z/2$}{pi1(SO(3)) = Z/2}}
$S^3$ is simply connected (Lee, Proposition~A.68(b), p.~614), hence so is
$SU(2)$ by Eq.~\eqref{quat:SU2-diffeo-S3}. So $\Phi$ is the universal
covering of $SO(3)$, and $SU(2)$ is the universal covering group of Lee's
Theorem~7.7 (p.~154).

\emph{An explicit non-contractible loop.} The rotations $R(2\theta)$ about the
$z$ axis, $0\le\theta\le\pi$, form a closed loop $\gamma$ in $SO(3)$ based at
$I$. By the example of \S\ref{quat:sec:covering}, its lift starting at $I$ is
$u(\theta)=\operatorname{diag}(e^{-i\theta},e^{i\theta})$, which ends at
$u(\pi)=-I$. By Lee's monodromy theorem (Proposition~A.77(c), p.~616),
path-homotopic paths have lifts with the same endpoint. The constant loop
lifts to the constant path at $I$, which ends at $I\ne-I$. So $\gamma$ is not
null-homotopic, and $\pi_1(SO(3))\ne1$. The loop $\gamma*\gamma$ (two full
turns) lifts to $u(\theta)$, $0\le\theta\le2\pi$, a closed loop in the simply
connected $SU(2)$. Contract it there and push the contraction down by $\Phi$:
$\gamma*\gamma$ is null-homotopic.

\emph{The whole group.} The covering automorphisms of $\Phi$ are $u\mapsto u$
and $u\mapsto-u$. A covering automorphism is determined by the image of one
point (unique lifting, Lee, Proposition~A.77(a)), and it must send $I$ into
$\Phi^{-1}(I)=\{\pm I\}$. Lee's Problem~21-19 (p.~562) quotes the fact that
for a universal covering the automorphism group is isomorphic to the
fundamental group of the base (proved in Lee's \emph{Introduction to
Topological Manifolds}, Ch.~12). Hence
\begin{equation}\label{quat:pi1-SO3}
 \pi_1\bigl(SO(3)\bigr)\cong\ker\Phi=\{\pm I\}\cong\mathbb Z/2 .
\end{equation}
This is the topological fact behind spin~$\tfrac12$. A continuous family of
rotations of $\mathbb R^3$ that starts and ends at $I$ lifts to $SU(2)$ and
ends at $I$ or at $-I$. A $2\pi$ turn ends at $-I$; a $4\pi$ turn ends at $I$.

\emph{Lie algebras do not see this.} $d\Phi$ identifies $\mathfrak{su}(2)$ with
$\mathfrak{so}(3)$. Since $SU(2)$ is simply connected, Lee's Theorem~20.19
(p.~530) lifts every Lie algebra homomorphism out of $\mathfrak{su}(2)$ to a
Lie group homomorphism out of $SU(2)$. The inverse $d\Phi^{-1}\colon
\mathfrak{so}(3)\to\mathfrak{su}(2)$ has no such lift out of $SO(3)$: a lift
$\Psi\colon SO(3)\to SU(2)$ with $\Psi_*=d\Phi^{-1}$ would give the
homomorphism $\Phi\circ\Psi$, whose induced Lie algebra map is the identity.
By naturality of $\exp$ (Lee, Proposition~20.8(g), p.~520),
$\Phi\circ\Psi(\exp X)=\exp X$ for every $X\in\mathfrak{so}(3)$, so
$\Phi\circ\Psi=\mathrm{Id}$ on the image of $\exp$, which contains a
neighborhood of $I$ (Proposition~20.8(f)). A neighborhood of $I$ generates the
connected $SO(3)$ (Lee, Proposition~7.14(c), p.~156), so
$\Phi\circ\Psi=\mathrm{Id}$ everywhere. Then
$\Psi\circ\gamma$ would be a closed lift of $\gamma$, contradicting the
monodromy argument above.

\section{Compactness, and what it buys}\label{quat:sec:compactness}
A topological space is compact if every open cover has a finite subcover
(Lee, p.~608). Compactness is a finiteness condition. A fact that holds near
each point, with its own neighborhood and its own constant, holds on finitely
many of those neighborhoods, so it holds everywhere with the worst of finitely
many constants. The same condition says that an infinite process cannot
escape. Suppose a sequence has no cluster point. Then each point has an open
neighborhood containing only finitely many terms, and these neighborhoods form
an open cover with no finite subcover.

\begin{proposition}\label{quat:prop:compact}
$SU(2)$ and $SO(3)$ are compact.
\end{proposition}
\begin{proof}
$SO(3)=\{g\in\mathbb R^{3\times3}: g^Tg=I,\ \det g=1\}$ is the common zero
set of the continuous functions $g\mapsto g^Tg-I$ and $g\mapsto\det g-1$, so
it is closed in $\mathbb R^9$. The columns of $g$ are unit vectors, so
$\|g\|_F^2=\operatorname{tr}(g^Tg)=\operatorname{tr}I=3$, and $SO(3)$ is
bounded. By Heine--Borel it is compact. In the same way, $SU(2)\approx S^3$ is
closed and bounded in $\mathbb R^4$, so it is compact. Alternatively,
$SO(3)=\Phi\bigl(SU(2)\bigr)$ is the continuous image of a compact space (Lee,
Proposition~A.45(a), p.~608). Hall gives the same argument for $O(n)$,
$SO(n)$, $U(n)$, $SU(n)$ \cite[\S1.3, p.~16]{BCHall2015}.
\end{proof}
Compactness was used once already. In \S\ref{quat:sec:covering}, the
continuous bijection $\mathbb RP^3\to SO(3)$ is a homeomorphism because its
domain is compact and its target is Hausdorff (Lee's closed map lemma,
Lemma~A.52, p.~610).

\subsection{No global three-parameter chart}
\begin{proposition}\label{quat:prop:no-global-chart}
No chart of $SO(3)$ has all of $SO(3)$ as its domain.
\end{proposition}
\begin{proof}
A chart is a homeomorphism $\varphi$ from an open subset of $SO(3)$ onto an
open subset $\varphi(SO(3))$ of $\mathbb R^3$. If its domain were all of
$SO(3)$, then $\varphi(SO(3))$ would be compact by
Proposition~\ref{quat:prop:compact}, hence closed in $\mathbb R^3$. A nonempty
subset of the connected $\mathbb R^3$ that is open and closed is
$\mathbb R^3$, which is not compact.
\end{proof}
Euler angles, rotation vectors, Rodrigues--Gibbs parameters and modified
Rodrigues parameters are charts, so each of them is singular or multivalued
somewhere. Gimbal lock is the boundary of the Euler-angle chart.

\subsection{No continuous choice of sign}
$\Phi$ is a local diffeomorphism everywhere, so four parameters with the
constraint $u\in SU(2)$ describe every rotation without singularity. The price
is the sign: $\Phi(-u)=\Phi(u)$.
\begin{proposition}\label{quat:prop:no-section}
There is no continuous map $s\colon SO(3)\to SU(2)$ with
$\Phi\circ s=\mathrm{Id}_{SO(3)}$.
\end{proposition}
\begin{proof}
Let $\gamma(\theta)=R(2\theta)$, $0\le\theta\le\pi$, be the loop of
\S\ref{quat:sec:smooth-cover}, with lift
$u(\theta)=\operatorname{diag}(e^{-i\theta},e^{i\theta})$ from $u(0)=I$ to
$u(\pi)=-I$. Suppose $s$ exists. Then $s\circ\gamma$ is continuous, closed
($s\circ\gamma(0)=s\circ\gamma(\pi)=s(I)$), and satisfies
$\Phi\circ(s\circ\gamma)=\gamma$, so it is a lift of $\gamma$ starting at
$s(I)\in\{\pm I\}$. By unique lifting (Lee, Proposition~A.77(a)), it equals
$u$ if $s(I)=I$, which ends at $-I$, and $-u$ if $s(I)=-I$, which ends at $I$.
Either way $s\circ\gamma$ is not closed, a contradiction.
\end{proof}
The proof uses only the restriction of $s$ to the loop $\gamma$. So every rule
that assigns one unit quaternion to each rotation is discontinuous somewhere
on every full turn about an axis. For example, the rule
``$\operatorname{tr}u\ge0$'' (scalar part $\ge0$) jumps from $u$ to $-u$ at
the rotations by $\pi$, where $\operatorname{tr}u(\theta)=2\cos\theta$ changes
sign. A continuous path $g(t)$ in $SO(3)$ does have a continuous lift $u(t)$,
unique once $u(0)$ is chosen (Lee, Proposition~A.77(a)(b)). A propagated
attitude can therefore be continuous in $SU(2)$, while the sign of a quaternion
built from a rotation matrix cannot be chosen continuously.

\subsection{Compactness and quantization}
\begin{remark}\label{quat:rem:compact-quantized}
For a compact group, averaging an inner product over the group, with the
invariant integral in place of a finite sum, makes every finite-dimensional
representation unitary, hence completely reducible (Hall, Theorem~4.28,
p.~92, proved this way). By the Peter--Weyl theorem, every irreducible unitary
representation of a compact Lie group is finite-dimensional, and the
irreducibles form a discrete family. For $SU(2)$ they are labeled by
$m=0,1,2,\ldots$, of dimension $m+1$ (Hall, Theorem~4.32, p.~96), that is,
$j=m/2\in\{0,\frac12,1,\ldots\}$. The integer-$j$ representations factor
through $SO(3)$; the half-integer ones do not (Hall, Proposition~4.35,
p.~101), because $-I\in\ker\Phi$ acts as $(-1)^{m}$. A connected noncompact
simple Lie group, such as the Lorentz group, has no nontrivial
finite-dimensional unitary representations, so its unitary representations of
physical interest are infinite-dimensional.

The same mechanism discretizes spectra. The Laplacian of a compact Riemannian
manifold without boundary has discrete eigenvalues of finite multiplicity,
because its
resolvent is a compact operator. On $S^2$ the eigenvalues of
$-\Delta_{S^2}$ are $\ell(\ell+1)$, the spectrum of $\mathbf L^2/\hbar^2$.
On the noncompact $\mathbb R^n$ the spectrum of $-\Delta$ is continuous. The
contrast between Fourier series on the compact $S^1$ (frequencies
$\mathbb Z$) and the Fourier transform on $\mathbb R$ (frequencies
$\mathbb R$) is the abelian case. A confining potential produces a discrete
spectrum through the same condition, a compact resolvent, without a compact
configuration space.
\end{remark}

\section{Dimension accounting}\label{quat:sec:dimensions}
\begin{center}
\small
\begin{tabular}{llll}
 space & ambient & conditions & dimension\\\hline
 $SO(3)$ & $9$ real entries & $g^Tg=I$: $6$; $\det g=1$: none & $3$\\
 $\mathfrak{so}(3)$ & $9$ & $X^T=-X$: $6$ & $3$\\
 $SU(2)$ & $8$ real & $U^\dagger U=1$: $4$; $\det U=1$: $1$ & $3$\\
 $\mathfrak{su}(2)$ & $8$ & $X^\dagger=-X$: $4$; $\operatorname{tr}X=0$: $1$ & $3$\\
 $S^3$ & $\mathbb R^4$ & $|z_1|^2+|z_2|^2=1$: $1$ & $3$\\
 $\ker\Phi$ & --- & $\{\pm I\}$, two points & $0$
\end{tabular}
\end{center}
$\Phi$ maps a $3$-dimensional group onto a $3$-dimensional group with a
$0$-dimensional kernel, and $d\Phi$ is an isomorphism of the
$3$-dimensional Lie algebras. The Lie algebras are isomorphic and the groups
are not: $d\Phi$ only describes $\Phi$ near the identity, and the two-point
kernel is global information that $\mathfrak{su}(2)\cong\mathfrak{so}(3)$ does
not detect.
